\chapter{EKSTENSI}

\section{Operator Normal secara Esensial}
\begin{defn} Misalkan $T$ suatu operator pada ruang Hilbert~$H$.  
 \index{spektrum!esensial}%
 \index{esensial!spektrum}%
 \index{sigmae@$\sigma_e(T)$ (spektrum esensial dari $T$)}%
\df{spektrum esensial} dari $T$, yang dinotasikan dengan $\sigma_e(T)$, adalah spektrum citra
$T$ dalam aljabar Calkin $\ofml Q(H)$; yaitu,
   \[ \sigma_e(T) = \sigma_{\ofml Q(H)}(\pi(T))\,. \]
\end{defn}

\begin{prop} Jika $T$ suatu operator pada ruang Hilbert $H$, maka
  \[ \sigma_e(T) = \{\lambda \in \C\colon T - \lambda I \notin \ofml F(H)\}\,. \]
\end{prop}

\begin{prop} Spektrum esensial dari operator swaadjoin $T$ pada ruang Hilbert adalah gabungan
titik-titik akumulasi spektrum $T$ dengan nilai-nilai eigen $T$ yang memiliki
multiplisitas tak hingga.  Anggota-anggota $\sigma(T) \setminus \sigma_e(T)$ adalah nilai-nilai eigen
terisolasi dengan multiplisitas hingga.
\end{prop}

\begin{proof} Lihat \cite{HigsonR:2000}, proposisi 2.2.2. \ns \end{proof}

Berikut dua teorema dari analisis fungsional klasik.

\begin{thm}[Weyl]\label{005121} Jika $S$ dan $T$ merupakan operator pada suatu ruang Hilbert
 \index{teorema Weyl}%
yang selisihnya kompak, maka spektrum keduanya sama, kecuali mungkin pada nilai-nilai eigen.
\end{thm}

\begin{thm}[Weyl--von Neumann]\label{005124} Misalkan $T$ suatu operator swaadjoin pada
 \index{teorema Weyl--von Neumann}%
ruang Hilbert terpisahkan~$H$. Untuk setiap $\epsilon > 0$ terdapat operator $D$ yang dapat
didiagonalkan sedemikian sehingga $T - D$ kompak dan $\norm{T - D} < \epsilon$.
\end{thm}

\begin{proof} Lihat \cite{Conway:2000}, 38.1; \cite{Davidson:1996}, akibat II.4.2; atau
\cite{HigsonR:2000}, 2.2.5. \ns
\end{proof}

\begin{defn} Misalkan $H$ dan $K$ ruang-ruang Hilbert. Operator $S \in \ofml B(H)$ dan $T \in \ofml
B(K)$ disebut
 \index{esensial!ekuivalen uniter}%
 \index{uniter!ekuivalensi!esensial}%
 \index{ekuivalensi!uniter esensial}%
\df{ekuivalen uniter secara esensial} (atau
 \index{kompalen}%
\df{kompalen}) jika terdapat pemetaan uniter $U\colon H \sto K$ sedemikian sehingga $S - U^*TU$
merupakan operator kompak pada~$H$. (Kita memperluas definisi~\ref{000161} dan~\ref{0001612} dengan
cara yang jelas: $U \in \ofml B(H,K)$ disebut
 \index{uniter!pemetaan linear terbatas}%
 \index{terbatas!pemetaan linear!uniter}%
\df{uniter} jika $U^*U = I_H$ dan $UU^* = I_K$; dan $S$ serta $T$ disebut
 \index{uniter!ekuivalensi!operator}%
 \index{ekuivalensi!uniter}%
\df{ekuivalen secara uniter} jika terdapat pemetaan uniter $U\colon H \sto K$ sedemikian sehingga $S =
U^*TU$.)
\end{defn}

\begin{prop}\label{005134} Operator-operator swaadjoin $S$ dan $T$ pada ruang-ruang Hilbert terpisahkan berdimensi tak hingga
ekuivalen uniter secara esensial jika dan hanya jika keduanya memiliki spektrum esensial yang sama.
\end{prop}

\begin{proof} Lihat \cite{HigsonR:2000}, proposisi 2.2.4.  \ns \end{proof}

\begin{defn} Suatu operator ruang Hilbert $T$ disebut
 \index{esensial!normal}%
 \index{normal!secara esensial}
 \index{operator!normal secara esensial}%
\df{normal secara esensial} jika
 \index{komutator}%
 \index{<brackets@$[T,T^*]$ (komutator dari $T$)}%
\emph{komutator}-nya $[T,T^*] := TT^* - T^*T$ kompak. Dengan kata lain: $T$ normal secara
esensial jika citranya $\pi(T)$ merupakan elemen normal dari aljabar Calkin. Operator $T$
disebut
 \index{esensial!swaadjoin}%
 \index{swaadjoin!secara esensial}
 \index{operator!swaadjoin secara esensial}%
\df{swaadjoin secara esensial} jika $T - T^*$ kompak; yakni, jika $\pi(T)$
swaadjoin dalam aljabar Calkin.
\end{defn}

\begin{exam} Operator geser unilateral $S$ (lihat contoh~\ref{C063526})
 \index{geser unilateral!normalitas esensial}%
normal secara esensial (tetapi tidak normal).
\end{exam}





















\section{Operator Toeplitz}
\begin{defn} Misalkan
 \index{zeta@$\zeta$ (fungsi identitas pada lingkaran satuan)}%
$\zeta\colon \T \sto \T\colon z \mapsto z$ merupakan fungsi identitas pada lingkaran satuan.
Maka $\{\zeta^n\colon n \in \Z\}$ merupakan basis ortonormal bagi ruang Hilbert $L_2(\T)$ yang terdiri atas
(kelas-kelas ekuivalensi dari) fungsi-fungsi yang terintegralkan kuadrat pada $\T$ terhadap ukuran panjang
busur (yang dinormalisasi secara sesuai). Kita menyatakan dengan $H^2$ subruang dari $L_2(\T)$ yang merupakan
rentang linear tertutup dari $\{\zeta^n\colon n \ge 0\}$ dan dengan $P_+$ proyeksi
(ortogonal) pada $L_2(\T)$ yang jangkauan
(sekaligus kodomain)-nya adalah~$H^2$. Ruang $H^2$ merupakan
contoh suatu
 \index{ruang Hardy!$H^2$}%
 \index{ruang!Hardy}%
 \index{H@$H^2$}%
\emph{ruang Hardy}. Untuk setiap $\phi \in L_\infty(\T)$ kita mendefinisikan pemetaan $T_\phi$ dari
ruang Hilbert $H^2$ ke dirinya sendiri dengan $T_\phi = P_+M_\phi$ (dengan $M_\phi$ operator
perkalian yang didefinisikan dalam contoh~\ref{C063527}). Operator demikian disebut
 \index{Toeplitz!operator}%
 \index{operator!Toeplitz}%
 \index{T@$T_\phi$ (operator Toeplitz dengan simbol~$\phi$)}%
\df{operator Toeplitz} dan $\phi$ disebut
 \index{simbol!operator Toeplitz}%
 \index{Toeplitz!operator!simbol}%
\df{simbol}-nya. Jelas bahwa $T_\phi$ merupakan operator pada $H^2$ dan $\norm{T_\phi} \le \norm
{\phi}_\infty$.
\end{defn}

\begin{exam} Operator Toeplitz $T_\zeta$
 \index{geser unilateral!sebagai operator Toeplitz}%
 \index{Toeplitz!operator!geser unilateral sebagai}%
bertindak pada vektor-vektor basis $\zeta^n$ dari $H^2$ menurut $T_\zeta(\zeta^n) = \zeta^{n+1}$, sehingga operator ini
ekuivalen secara uniter dengan geser unilateral $S$.
\end{exam}

\begin{prop}\label{005216} Pemetaan $T\colon L_\infty(\T) \sto \ofml B(H^2)\colon \phi
\mapsto T_\phi$ bersifat positif, linear, dan melestarikan involusi.
\end{prop}

\begin{exam}\label{005217} Operator Toeplitz $T_\zeta$ dan $T_{\conj\zeta}$ menunjukkan bahwa
pemetaan $T$ dalam proposisi sebelumnya \emph{bukan} representasi dari $L_\infty(\T)$
pada~$\ofml B(H^2)$. (Bandingkan dengan contoh~\ref{002802}.)
\end{exam}

\begin{notn} Misalkan $H^\infty := H^2 \cap L_\infty(\T)$.
 \index{H@$H^\infty$}%
 \index{ruang Hardy!$H^\infty$}%
 \index{ruang!Hardy}%
Ini merupakan contoh lain dari suatu \emph{ruang Hardy}.
\end{notn}

\begin{prop} Suatu fungsi terbatas secara esensial $\phi$ pada lingkaran satuan termasuk dalam
$H^\infty$ jika dan hanya jika operator perkalian $M_\phi$ memetakan $H^2$ ke dalam~$H^2$.
\end{prop}

Meskipun (seperti telah kita lihat dalam contoh~\ref{005217}) pemetaan $T$ yang didefinisikan dalam
proposisi~\ref{005216} pada umumnya tidak multiplikatif, pemetaan ini multiplikatif untuk suatu
kelas fungsi yang luas.

\begin{prop} Jika $\phi \in L_\infty(\T)$ dan $\psi \in H^\infty$, maka $T_{\phi\psi} =
T_\phi T_\psi$ dan $T_{\conj\psi \phi} = T_{\conj\psi} T_\phi$.
\end{prop}

\begin{prop} Jika operator Toeplitz $T_\phi$ dengan simbol $\phi \in L_\infty(\T)$
invertibel, maka fungsi $\phi$ invertibel.
\end{prop}

\begin{proof} Lihat \cite{Douglas:1972}, proposisi 7.6. \ns \end{proof}

\begin{thm}[Inklusi Spektral Hartman--Wintner] Jika $\phi \in L_\infty(\T)$,
 \index{teorema inklusi spektral Hartman--Wintner}%
maka $\sigma(\phi) \subseteq \sigma(T_\phi)$ dan $\rho(T_\phi) = \norm{T_\phi} =
\norm{\phi}_\infty$.
\end{thm}

\begin{proof} Lihat \cite{Douglas:1972}, akibat 7.7 atau \cite{Murphy:1990}, teorema 3.5.7. \ns
\end{proof}

\begin{cor} Pemetaan $T\colon L_\infty(\T) \sto \ofml B(H^2)$ yang didefinisikan dalam
proposisi~\ref{005216} merupakan suatu isometri.
\end{cor}

\begin{prop} Suatu operator Toeplitz $T_\phi$ dengan simbol $\phi \in L_\infty(\T)$ kompak jika
dan hanya jika $\phi = \vc 0$.
\end{prop}

\begin{proof} Lihat \cite{Murphy:1990}, teorema 3.5.8.  \ns \end{proof}

\begin{prop} Jika $\phi \in \fml C(\T)$ dan $\psi \in L_\infty(\T)$, maka
semikomutator $T_\phi T_\psi - T_{\phi\psi}$ dan $T_\psi T_\phi - T_{\psi\phi}$ bersifat
kompak.
\end{prop}

\begin{proof} Lihat \cite{Arveson:2002}, proposisi 4.3.1; \cite{Davidson:1996}, akibat
V.1.4; atau \cite{Murphy:1990}, lema 3.5.9. \ns
\end{proof}

\begin{cor} Setiap operator Toeplitz dengan simbol kontinu bersifat normal secara esensial.
\end{cor}

\begin{defn} Andaikan $H$ ruang Hilbert terpisahkan dengan basis $\{e_0, e_1, e_2,
\dots\}$ dan $T$ suatu operator pada $H$ yang representasi matriks (tak hingga)-nya adalah
$[t_{ij}]$. Jika entri-entri dalam matriks ini hanya bergantung pada selisih indeks $i
- j$ (yakni, jika setiap diagonal yang sejajar dengan diagonal utama merupakan barisan konstan),
maka matriks tersebut disebut
 \index{Toeplitz!matriks}%
 \index{matriks!Toeplitz}%
\df{matriks Toeplitz}.
\end{defn}

\begin{prop} Misalkan $T$ suatu operator pada ruang Hilbert $H^2$. Representasi matriks
$T$ terhadap basis biasa $\{\zeta^n\colon n \ge 0\}$ merupakan matriks Toeplitz jika
dan hanya jika $S^*TS = T$ (dengan $S$ geser unilateral).
\end{prop}

\begin{proof} Lihat \cite{Arveson:2002}, proposisi 4.2.3.  \ns \end{proof}

\begin{prop} Jika $T_\phi$ suatu operator Toeplitz dengan simbol $\phi \in L_\infty(\T)$, maka
$S^*T_\phi S = T_\phi$. Sebaliknya, jika $R$ suatu operator pada ruang Hilbert $H^2$
sedemikian sehingga $S^*RS = R$, maka terdapat $\phi \in L_\infty(\T)$ tunggal sedemikian sehingga $R =
T_\phi$.
\end{prop}

\begin{proof} Lihat \cite{Arveson:2002}, teorema 4.2.4.  \ns \end{proof}

\begin{defn} 
 \index{Toeplitz!aljabar}%
 \index{aljabar!Toeplitz}%
 \index{T@$\ofml T$ (aljabar Toeplitz)}%
\df{aljabar Toeplitz} $\ofml T$ adalah subaljabar-$C^*$ dari $\ofml B(H^2)$ yang dibangkitkan oleh
operator geser unilateral; yaitu, $\ofml T = C^*(S)$.
\end{defn}

\begin{prop} Himpunan $\ofml K(H^2)$ dari operator-operator kompak pada $H^2$ merupakan ideal dalam aljabar
Toeplitz~$\ofml T$.
\end{prop}

\begin{prop} Aljabar Toeplitz terdiri atas semua perturbasi kompak dari operator Toeplitz
dengan simbol kontinu. Yaitu,
   \[ \ofml T = \{ T_\phi + K\colon \phi \in \fml C(\T) \text{ dan } K \in \ofml K(H^2)\} \,. \]
Lebih lanjut, jika $T_\phi + K = T_\psi + L$ dengan $\phi$, $\psi \in \fml C(\T)$ dan $K$, $L
\in \ofml K(H^2)$, maka $\phi = \psi$ dan $K = L$.
\end{prop}

\begin{proof} Lihat \cite{Arveson:2002}, teorema 4.3.2 atau \cite{Davidson:1996}, teorema V.1.5.
\ns \end{proof}

\begin{prop}\label{005251} Pemetaan $\pi \circ T\colon \fml C(\T) \sto \ofml Q(H^2)\colon \phi
\mapsto \pi(T_\phi)$ merupakan monomorfisme-$*\,$ beridentitas. Jadi pemetaan $\alpha\colon \fml C(\T)
\sto \ofml T/\ofml K(H^2)\colon \phi \mapsto \pi(T_\phi)$ menghasilkan isomorfisme
antara aljabar-$C^*$ $\fml C(\T)$ dan $\ofml T/\ofml K(H^2)$.
\end{prop}

\begin{proof} Lihat \cite{HigsonR:2000}, proposisi 2.3.3 atau \cite{Murphy:1990},
teorema 3.5.11.  \ns
\end{proof}

\begin{cor} Jika $\phi \in \fml C(\T)$, maka $\sigma_e(T_\phi) = \ran\phi$.
\end{cor}

\begin{prop}\label{005253} Barisan
   \[ \vc 0 \to \ofml K(H^2) \to \ofml T \to^\beta \fml C(\T) \to \vc 0 \]
bersifat eksak. Barisan ini tidak terbelah.
\end{prop}

\begin{proof} Pemetaan $\beta$ didefinisikan oleh $\beta(R) := \alpha^{-1}\bigl(\pi(R)\bigr)$ untuk
setiap $R \in \ofml T$ (dengan $\alpha$ isomorfisme yang didefinisikan dalam
proposisi sebelumnya~\ref{005251}). Lihat \cite{Arveson:2002}, catatan 4.3.3; \cite{Davidson:1996},
teorema V.1.5; atau \cite{HigsonR:2000}, halaman 35. \ns
\end{proof}

\begin{defn}\label{0052531} Barisan eksak pendek dalam proposisi sebelumnya~\ref{005253}
adalah
 \index{Toeplitz!ekstensi}%
 \index{ekstensi!Toeplitz}%
\df{ekstensi Toeplitz} dari $\fml C(\T)$ oleh~$\ofml K(H^2)$.
\end{defn}

\begin{rem} Suatu versi diagram ekstensi Toeplitz yang sering muncul kira-kira berbentuk
seperti
     \[ \vc 0 \to \ofml K(H^2) \to \ofml T \two/->`<-/^\beta_T \fml C(\T) \to \vc 0 \tag{1} \]
dengan $\beta$ seperti dalam~\ref{005253} dan $T$ adalah pemetaan $\phi \mapsto T_\phi$. Diagram ini
dapat disalahtafsirkan. Diagram tersebut mungkin memberi kesan kepada pembaca yang kurang waspada bahwa ini adalah ekstensi
terbelah, terutama karena memang benar bahwa $\beta \circ T = I_{\fml
C(\T)}$. Masalahnya, tentu saja, bahwa diagram ini bukan diagram dalam kategori
$\cat{CSA}$ dari aljabar-$C^*$ dan homomorfisme-$*\,$. Kita telah melihat dalam
contoh~\ref{005217} bahwa pemetaan $T$ bukan homomorfisme-$*\,$ karena tidak
selalu melestarikan perkalian. Elemen-elemen invertibel dalam $\fml C(\T)$ belum tentu terangkat menjadi
elemen-elemen invertibel dalam aljabar Toeplitz~$\ofml T$; jadi epimorfisme-$*\,$ $\beta$
tidak memiliki invers kanan dalam kategori~$\cat{CSA}$.

Sebagian penulis menangani masalah ini dengan mengatakan bahwa barisan (1) bersifat semiterbelah (lihat, misalnya,
Arveson~\cite{Arveson:2002}, halaman 112). Penulis lain meminjam istilah dari teori kategori,
yang menggunakan kata ``penampang'' untuk berarti ``memiliki invers kanan''.
Davidson~\cite{Davidson:1996}, misalnya, pada halaman 134 menyebut pemetaan $T$ sebagai
``penampang kontinu'', sedangkan Douglas~\cite{Douglas:1972}, pada halaman 179 menyebutnya
``penampang silang isometrik''. Meskipun memang benar bahwa $\beta$ memiliki $T$ sebagai invers
kanan dalam kategori $\cat{SET}$ dari himpunan dan pemetaan, bahkan juga dalam kategori ruang Banach
dan pemetaan linear terbatas, pemetaan tersebut tidak memiliki invers kanan dalam $\cat{CSA}$.
\end{rem}

\begin{prop} Jika $\phi$ suatu fungsi dalam $\fml C(\T)$, maka $T_\phi$ Fredholm jika dan hanya jika
fungsi itu tidak pernah bernilai nol.
\end{prop}

\begin{proof} Lihat \cite{Arveson:2002}, halaman 112, akibat 1; \cite{Davidson:1996}, teorema
V.1.6; \cite{Douglas:1972}, teorema 7.26; atau \cite{Murphy:1990}, akibat 3.5.12. \ns
\end{proof}

\begin{prop}\label{005271} Jika $\phi$ elemen invertibel dari $\fml C(\T)$, maka terdapat
bilangan bulat $n$ tunggal sedemikian sehingga $\phi = \zeta^n \exp \psi$ untuk suatu $\psi \in \fml
C(\T)$.
\end{prop}

\begin{proof} Lihat \cite{Arveson:2002}, proposisi 4.4.1 dan 4.4.2 atau \cite{Murphy:1990},
lema 3.5.14. \ns
\end{proof}

\begin{defn} Bilangan bulat yang keberadaannya dinyatakan dalam proposisi sebelumnya~\ref{005271}
adalah
 \index{bilangan lilit}%
 \index{w@$w(\phi)$ (bilangan lilit dari~$\phi$)}%
\df{bilangan lilit} dari fungsi invertibel $\phi \in \fml C(\T)$. Bilangan ini dinotasikan
dengan~$w(\phi)$.
\end{defn}

\begin{thm}[Teorema indeks Toeplitz]\label{005274}
 \index{Toeplitz!teorema indeks}%
Jika $T_\phi$ suatu operator Toeplitz dengan simbol kontinu yang tidak pernah bernilai nol, maka operator ini
merupakan operator Fredholm dan
   \[ \ind(T_\phi) = -w(\phi)\,. \]
\end{thm}

\begin{proof} Bukti elementer dapat ditemukan dalam \cite{Arveson:2002}, teorema 4.4.3, dan
\cite{Murphy:1990}. Bagi pembaca yang memiliki sedikit latar belakang tentang homotopi kurva,
proposisi~\ref{005271}, yang mengantar ke definisi \emph{bilangan lilit} di atas,
dapat dilewati. Merupakan fakta elementer dalam teori homotopi bahwa grup fundamental
$\pi_1(\C \setminus \{0\})$ dari bidang kompleks berlubang bersifat siklik tak hingga. Terdapat
isomorfisme $\tau$ dari $\pi_1(\C \setminus \{0\})$ ke $\Z$ yang mengaitkan bilangan bulat $+1$
dengan (kelas ekuivalensi yang memuat) fungsi $\zeta$. Hal ini memungkinkan kita untuk
mengaitkan setiap anggota invertibel $\phi$ dari $\fml C(\T)$ dengan bilangan bulat yang bersesuaian
di bawah $\tau$ dengan kelas ekuivalensinya. Bilangan bulat ini kita sebut
 \index{bilangan lilit}%
 \index{fundamental!grup}%
 \index{grup!fundamental}%
\emph{bilangan lilit} dari~$\phi$. Definisi demikian memungkinkan pemberian bukti yang lebih singkat
dan lebih anggun bagi \emph{teorema indeks Toeplitz}. Untuk pembahasan dengan pendekatan ini,
lihat \cite{Douglas:1972}, teorema 7.26, atau \cite{HigsonR:2000}, teorema 2.3.2. Untuk
latar belakang mengenai grup fundamental, lihat \cite{Massey:1967}, bab dua;
\cite{Wallace:1970}, lampiran A; atau \cite{Willard:1968}, bagian 32 dan 33. \ns
\end{proof}

\begin{thm}[Dekomposisi Wold]
 \index{dekomposisi Wold}%
 \index{dekomposisi!Wold}%
Setiap
 \index{sejati!isometri}%
 \index{isometri!sejati}%
isometri sejati (yakni, nonuniter) pada suatu ruang Hilbert merupakan jumlah langsung salinan-salinan
geser unilateral, atau merupakan jumlah langsung suatu operator uniter bersama salinan-salinan
geser tersebut.
\end{thm}

\begin{proof} Lihat \cite{Davidson:1996}, teorema V.2.1; \cite{Halmos:1982}, soal (dan
solusi) 149; atau \cite{Murphy:1990}, teorema 3.5.17. \ns
\end{proof}

\begin{thm}[Coburn]
 \index{teorema Coburn}%
Misalkan $v$ suatu isometri dalam aljabar-$C^*$ beridentitas~$A$.
Maka
terdapat
homomorfisme-$*\,$
beridentitas
$\tau$ tunggal dari aljabar Toeplitz $\ofml T$ ke $A$ sedemikian sehingga
$\tau(T_\zeta) = v$. Lebih lanjut, jika $vv^* \ne \vc 1$, maka $\tau$ merupakan suatu isometri.
\end{thm}

\begin{proof} Lihat \cite{Murphy:1990}, teorema 3.5.18, atau \cite{Davidson:1996},
teorema V.2.2. \ns
\end{proof}
%
%%
 
  
 











\section{Penjumlahan Ekstensi}

 \index{konvensi!ruang Hilbert separabel dan berdimensi tak hingga (mulai bagian Penjumlahan Ekstensi)}%
\begin{center}
 \fbox{\parbox{5.5in}{\textbf{Mulai sekarang semua ruang Hilbert akan separabel dan berdimensi
 tak hingga.}}}
\end{center}

\begin{prop} Suatu operator ruang Hilbert $T$ bersifat swaadjoin secara esensial jika dan hanya jika
operator itu merupakan perturbasi kompak dari suatu operator swaadjoin.
\end{prop}

\begin{prop}\label{005414} Dua operator ruang Hilbert $T_1$ dan $T_2$ yang swaadjoin secara esensial
ekuivalen uniter secara esensial jika dan hanya jika keduanya memiliki spektrum esensial yang
sama.
\end{prop}

Hasil yang analog tidak berlaku bagi operator normal secara esensial.

\begin{exam} Operator Toeplitz $T_\zeta$ dan $T_{\zeta^2}$ bersifat normal secara esensial dengan
spektrum esensial yang sama, tetapi keduanya tidak ekuivalen uniter secara esensial.
\end{exam}

\begin{defn} Sekarang kita membatasi perhatian pada suatu kelas khusus ekstensi. Jika $H$ merupakan
ruang Hilbert dan $A$ merupakan aljabar-$C^*$, kita mengatakan bahwa $(\ofml E,\phi)$ adalah suatu
 \index{ekstensi!$\ofml K(H)$ oleh $A$}%
\df{ekstensi $\ofml K = \ofml K(H)$ oleh $A$} jika $\ofml E$ merupakan subaljabar-$C^*$ beridentitas
dari $\ofml B(H)$ yang memuat $\ofml K$ dan $\phi$ merupakan homomorfisme-$*\,$ beridentitas sedemikian sehingga
barisan
 \begin{equation}\label{005431i}
    \vc 0 \to \ofml K \to^\iota \ofml E \to^\phi A \to \vc 0
 \end{equation}
(dengan $\iota$ pemetaan inklusi) eksak. Bahkan, kita hampir secara eksklusif akan
memperhatikan kasus $A = \fml C(X)$ untuk suatu ruang metrik kompak~$X$. Kita
mengatakan bahwa dua ekstensi $(\ofml E,\phi)$ dan $(\ofml E\,',\phi\,')$
 \index{ekuivalensi!ekstensi}%
 \index{ekstensi!ekuivalensi}%
\df{ekuivalen} jika terdapat suatu isomorfisme $\psi\colon \ofml E \sto \ofml E\,'$ yang
membuat diagram berikut komutatif.
 \begin{equation}\label{005431ii}
   \xymatrix{
     \vc 0\ar[r] & \ofml K\ar[r]^\iota\ar[d]_{\psi|_{\ofml K}} & \ofml E\ar[r]^\phi\ar[d]^\psi
                 & A\ar[r]\ar@{=}[d] & \vc 0 \\
     \vc 0\ar[r] & \ofml K\ar[r]^\iota & \ofml E\,'\ar[r]^{\phi\,'} & A\ar[r] & \vc 0
 }\end{equation}
Perhatikan bahwa definisi ini sedikit berbeda dari \emph{ekuivalensi kuat}
pada~\ref{0015151}. Keluarga semua kelas ekuivalensi ekstensi semacam ini kita lambangkan
 \index{ext@$\ext A$, $\ext X$!kelas ekuivalensi ekstensi}%
dengan~$\ext A$. Jika $A = \fml C(X)$, kita menulis $\ext X$, bukan~$\ext \fml C(X)$.
\end{defn}

\begin{defn} Jika $U\colon H_1 \sto H_2$ merupakan pemetaan uniter di antara ruang-ruang Hilbert, maka pemetaan
   \[ \ad_U\colon \ofml B(H_1) \sto \ofml B(H_2)\colon T \mapsto UTU^* \]
disebut
 \index{konjugasi}%
 \index{adU@$\ad_U$ (konjugasi oleh $U$)}%
\df{konjugasi} oleh~$U$.
\end{defn}

Jelas bahwa $\ad_U$ merupakan isomorfisme di antara aljabar-$C^*$ $\ofml B(H_1)$ dan
$\ofml B(H_2)$. Khususnya, jika $U$ merupakan operator uniter pada $H_1$, maka $\ad_U$ merupakan
automorfisme dari $\ofml K(H_1)$ maupun $\ofml B(H_1)$. Lebih lanjut, konjugasi merupakan
satu-satunya automorfisme aljabar-$C^*$~$\ofml K(H)$.

\begin{prop}\label{005436} Jika $H$ merupakan ruang Hilbert dan $\phi\colon \ofml K(H) \sto \ofml K(H)$
merupakan suatu automorfisme, maka $\phi = \ad_U$ untuk suatu operator uniter $U$ pada~$H$.
\end{prop}

\begin{proof} Lihat \cite{Davidson:1996}, lema V.6.1.  \ns \end{proof}

\begin{prop} Jika $H$ merupakan ruang Hilbert dan $A$ merupakan aljabar-$C^*$, maka ekstensi
$(\ofml E,\phi)$ dan $(\ofml E\,',\phi\,')$ dalam $\ext A$ ekuivalen jika dan hanya jika
terdapat suatu operator uniter $U$ dalam $\ofml B(H)$ sedemikian sehingga $\ofml E\,' = U \ofml E
U^*$ dan $\phi = \phi\,' \ad_U$.
\end{prop}

\begin{exam} Andaikan $T$ merupakan operator normal secara esensial pada suatu ruang Hilbert~$H$. Misalkan
 \index{E@$\ofml E_T$ (aljabar-$C^*$ yang dibangkitkan oleh $T$ dan $\ofml K(H)$)}%
$\ofml E_T$ merupakan aljabar-$C^*$ beridentitas yang dibangkitkan oleh $T$ dan~$\ofml K(H)$. Karena
$\pi(T)$ merupakan elemen normal aljabar Calkin, aljabar-$C^*$ beridentitas $\ofml
E_T/\ofml K(H)$ yang dibangkitkannya bersifat komutatif. Jadi \emph{teorema spektral
abstrak}~\ref{00144} memberi kita suatu isomorfisme aljabar-$C^*$ $\Psi\colon \fml C(\sigma_e(T))
\sto \ofml E_T/\ofml K(H)$. Misalkan $\sbsb{\phi}T = \Psi^{-1} \circ\pi\bigr|_{\ofml E_T}$.
Maka barisan
  \[ \vc 0 \to \ofml K(H) \to^\iota \ofml E_T \to^{\sbsb{\phi}T} \fml C(\sigma_e(T)) \to \vc 0 \]
eksak. Ini adalah
 \index{ekstensi!yang ditentukan oleh suatu operator normal secara esensial}%
\df{ekstensi $\ofml K(H)$ yang ditentukan oleh~$T$}.
\end{exam}

\begin{prop} Misalkan $T$ dan $T\,'$ operator normal secara esensial pada suatu ruang Hilbert~$H$. Kedua
operator itu ekuivalen uniter secara esensial jika dan hanya jika ekstensi yang
ditentukannya ekuivalen.
\end{prop}

\begin{prop} Jika $\ofml E$ merupakan aljabar-$C^*$ sedemikian sehingga $\ofml K(H) \subseteq \ofml E \subseteq
\ofml B(H)$ untuk suatu ruang Hilbert~$H$, $X$ merupakan subhimpunan kompak tak kosong dari $\C$, dan
$(\ofml E,\phi)$ merupakan ekstensi $\ofml K(H)$ oleh $\fml C(X)$, maka setiap elemen
dari $\ofml E$ bersifat normal secara esensial.
\end{prop}

\begin{defn} Jika $\phi_1\colon A_1 \sto B$ dan $\phi_2\colon A_2 \sto B$ merupakan
homomorfisme-$*\,$ beridentitas di antara aljabar-$C^*$, maka suatu
 \index{tarik balik}%
\df{tarik balik $A_1$ dan $A_2$ sepanjang $\phi_1$ dan $\phi_2$}, yang dilambangkan dengan
$(P,\pi_1,\pi_2)$, adalah aljabar-$C^*$ $P$ beserta sepasang
homomorfisme-$*\,$ beridentitas $\pi_1\colon P \sto A_1$ dan $\pi_2\colon P \sto A_2$ yang memenuhi
dua syarat berikut:
 \begin{enumerate}
    \item[(i)] diagram
        \[ \xy
          \Square[P`A_2`A_1`B;\pi_2`\pi_1`\phi_2`\phi_1]
      \endxy \]
bersifat komutatif; dan
    \item[(ii)] jika $\rho_1\colon Q \sto A_1$ dan $\rho_2\colon Q \sto A_2$ merupakan
homomorfisme-$*\,$ beridentitas dari aljabar-$C^*$ sedemikian sehingga diagram
        \[ \xy
          \Square[Q`A_2`A_1`B;\rho_2`\rho_1`\phi_2`\phi_1]
      \endxy \]
bersifat komutatif, maka terdapat suatu homomorfisme-$*\,$ beridentitas tunggal $\Psi\colon Q \sto P$
yang membuat diagram
       \[  \xy
         \pullback|brrb|<800,800>[P`A_2`A_1`B;\pi_2`\pi_1`\phi_2`\phi_1]%
         |amb|/{>}`{-->}`{>}/<500,500>[Q;\rho_2`\Psi`\rho_1]
       \endxy \]
komutatif.
 \end{enumerate}
\end{defn}

\begin{prop} Misalkan $H$ ruang Hilbert, $A$ aljabar-$C^*$ beridentitas, dan $\tau\colon
A \sto \ofml Q(H)$ monomorfisme-$*\,$ beridentitas. Maka, hingga isomorfisme,
terdapat tarik balik tunggal $(\ofml E,\pi_1,\pi_2)$ dari $A$ dan $\ofml B(H)$ sepanjang $\tau$ dan
$\pi$ sedemikian sehingga $(\ofml E,\pi_2)$ merupakan ekstensi $\ofml K(H)$ oleh $A$ yang membuat
diagram berikut komutatif.
 \begin{equation}\label{005464i}
   \xymatrix{
     \vc 0\ar[r] & \ofml K(H)\ar[r]^\iota\ar@{=}[d] & \ofml E\ar[r]^{\pi_2}\ar[d]^{\pi_1}
                 & A\ar[r]\ar[d]^\tau & \vc 0 \\
     \vc 0\ar[r] & \ofml K(H)\ar[r] & \ofml B(H)\ar[r]^{\pi} & \ofml Q(H)\ar[r] & \vc 0
 }\end{equation}
\end{prop}

\begin{proof}(Sketsa.) Misalkan $\ofml E = \{T \oplus a \in \ofml B(H) \oplus A\colon \tau(a) =
\pi(T)\}$, $\pi_1\colon \ofml E \sto \ofml B(H)\colon T \oplus a \mapsto T$, dan
$\pi_2\colon \ofml E \sto A\colon T \oplus a \mapsto a$. Ketunggalan (setiap
tarik balik) dibuktikan memakai ``omong kosong abstrak'' yang biasa.
\end{proof}

\begin{prop}\label{005466} Misalkan $H$ ruang Hilbert, $A$ aljabar-$C^*$ beridentitas, dan
$(\ofml E,\phi)$ suatu ekstensi $\ofml K(H)$ oleh~$A$. Maka terdapat suatu
monomorfisme-$*\,$ tunggal $\tau \colon A \sto \ofml Q(H)$ yang membuat
diagram~\eqref{005464i} komutatif.
\end{prop}

\begin{defn} Misalkan $H$ ruang Hilbert dan $A$ aljabar-$C^*$ beridentitas. Untuk $j=1,2$, misalkan
$\tau_j\colon A \sto \ofml Q(H)$ suatu monomorfisme-$*\,$ beridentitas. Kedua pemetaan ini
 \index{uniter!ekuivalensi!monomorfisme-$*\,$}%
 \index{ekuivalensi!uniter!monomorfisme-$*\,$}%
\df{ekuivalen secara uniter} jika terdapat suatu operator uniter $U$ pada $H$ sedemikian sehingga $\tau_2 =
\ad_U \tau_1$. Ekuivalensi uniter monomorfisme-$*\,$ beridentitas tentu merupakan suatu
relasi ekuivalensi. Kelas ekuivalensi yang memuat $\tau_1$ dilambangkan
dengan~$[\tau_1]$.
\end{defn}

\begin{prop} Misalkan $H$ ruang Hilbert dan $A$ aljabar-$C^*$ beridentitas. Dua ekstensi
$\ofml K(H)$ oleh $A$ ekuivalen jika dan hanya jika monomorfisme-$*\,$ beridentitas yang
bersesuaian dengannya (lihat~\ref{005466}) ekuivalen secara uniter.
\end{prop}

\begin{cor} Jika $H$ ruang Hilbert dan $A$ aljabar-$C^*$ beridentitas, terdapat suatu korespondensi
satu-satu di antara kelas-kelas ekuivalensi ekstensi $\ofml K(H)$ oleh $A$ dan
kelas-kelas ekuivalensi uniter monomorfisme-$*\,$ beridentitas dari $A$ ke aljabar
Calkin~$\ofml Q(H)$.
\end{cor}

\begin{conv} Berdasarkan akibat sebelumnya, kita akan memandang anggota $\ext A$ (atau
$\ext X$) sebagai
 \index{konvensi!E@$\ext A$ terdiri atas ekstensi atau morfisme}%
kelas ekuivalensi ekstensi atau kelas ekuivalensi uniter monomorfisme-$*\,$,
mana pun yang paling sesuai pada saat itu.
\end{conv}

\begin{prop}\label{0054702} Setiap ruang Hilbert $H$ (yang separabel dan berdimensi tak hingga)
isomorfik secara isometrik dengan $H \oplus H$. Jadi aljabar-$C^*$ $\ofml B(H)$ dan $\ofml
B(H \oplus H)$ isomorfik.
\end{prop}

\begin{defn}\label{0054715} Misalkan $\tau_1$, $\tau_2 \colon A \sto \ofml Q(H)$ merupakan
monomorfisme-$*\,$ beridentitas (dengan $A$ aljabar-$C^*$ dan $H$ ruang Hilbert). Kita mendefinisikan
suatu monomorfisme-$*\,$ beridentitas $\tau_1 \oplus \tau_2 \colon A \sto \ofml Q(H)$ melalui
  \[ (\tau_1 \oplus \tau_2)(a) := \rho\bigl(\tau_1(a) \oplus \tau_2(a)\bigr) \]
untuk setiap $a \in A$, dengan (seperti dalam Douglas\cite{Douglas:1980}) $\nu$ isomorfisme yang
ditetapkan dalam~\ref{0054702} dan $\rho$ pemetaan yang membuat diagram berikut
komutatif.
 \begin{equation}\label{0054715i}
  \xymatrix{
    \ofml B(H) \oplus \ofml B(H)\ar[r]\ar[dd]_{\pi \oplus \pi} & \ofml B(H \oplus H)\ar[r]^-\nu
                  & \ofml B(H)\ar[dd]^\pi   \\ & & & \\
    \ofml Q(H) \oplus \ofml Q(H)\ar[rr]_-\rho && \ofml Q(H)
 }\end{equation}
 Kemudian kita mendefinisikan operasi penjumlahan yang jelas pada $\ext A$:
   \[ [\tau_1] + [\tau_2] := [\tau_1 \oplus \tau_2] \]
untuk $[\tau_1]$, $[\tau_2] \in \ext A$.
\end{defn}

\begin{prop} Operasi penjumlahan (yang diberikan dalam~\ref{0054715}) pada $\ext A$ terdefinisi dengan baik
dan di bawah operasi ini $\ext A$ menjadi semigrup komutatif.
\end{prop}

\begin{defn}\label{0054811} Misalkan $A$ aljabar-$C^*$ beridentitas dan $\mathbf r\colon A \sto
\ofml B(H)$ representasi tak terdegenerasi dari $A$ pada suatu ruang Hilbert~$H$. Misalkan $P$
proyeksi ruang Hilbert dan $M$ jangkauan~$P$. Andaikan $P\mathbf r(a) -
\mathbf r(a)P \in \ofml K(H)$ untuk setiap $a~\in~A$. Lambangkan dengan $P_M$ proyeksi $P$
yang himpunan kodomainnya sama dengan jangkauannya; yaitu, $P_M\colon H \sto M \colon x \mapsto
Px$. Maka untuk setiap $a \in A$ kita mendefinisikan
 \index{operator!Toeplitz!abstrak}%
 \index{Toeplitz!operator!abstrak}%
 \index{abstrak!operator!Toeplitz}%
\df{operator Toeplitz abstrak} $T_a \in \ofml B(M)$
 \index{simbol!operator Toeplitz abstrak}%
 \index{T@$T_a$ (operator Toeplitz abstrak)}%
\df{bersimbol $a$ yang dikaitkan dengan} pasangan~$(\mathbf r,P)$ melalui $T_a = P_M\mathbf
r(a)\bigr|_M$.
\end{defn}

\begin{defn} Gunakan notasi seperti dalam definisi sebelumnya~\ref{0054811}. Kemudian kita mendefinisikan
 \index{ekstensi!Toeplitz!abstrak}%
 \index{Toeplitz!ekstensi!abstrak}%
 \index{abstrak!ekstensi!Toeplitz}%
\df{ekstensi Toeplitz abstrak $\tau_P$ yang dikaitkan dengan pasangan $(\mathbf r,P)$} melalui
   \[ \tau_P\colon A \sto \ofml Q(M)\colon a \mapsto \pi(T_a)\,. \]
\end{defn}

Perhatikan bahwa (monomorfisme-$*\,$ beridentitas yang dikaitkan dengan) ekstensi Toeplitz konkret
yang didefinisikan dalam~\ref{0052531} merupakan contoh ekstensi Toeplitz abstrak, dan
juga bahwa secara umum ekstensi Toeplitz abstrak tidak harus injektif.

\begin{defn} Misalkan $A$ aljabar-$C^*$ beridentitas dan $H$ ruang Hilbert. Dalam semangat
\cite{HigsonR:2000}, definisi 2.7.6, kita mengatakan bahwa suatu monomorfisme-$*\,$ beridentitas
$\tau\colon A \sto \ofml Q(H)$ bersifat
 \index{semiterbelah}%
\df{semiterbelah} jika terdapat suatu monomorfisme-$*\,$ beridentitas $\tau\,'\colon A \sto \ofml
Q(H)$ sedemikian sehingga $\tau \oplus \tau'$ terbelah.
\end{defn}

\begin{prop} Andaikan $A$ aljabar-$C^*$ beridentitas dan $H$ ruang Hilbert. Maka suatu
monomorfisme-$*\,$ beridentitas $\tau\colon A \sto \ofml Q(H)$ semiterbelah jika dan hanya jika
ekuivalen secara uniter dengan suatu ekstensi Toeplitz abstrak.
\end{prop}

\begin{proof} Lihat \cite{HigsonR:2000}, proposisi 2.7.10.  \ns \end{proof}

















\vskip .2 in


%\section{Hasil Kali Tensor Aljabar-$C^*$}


\vskip .2 in


















\section{Pemetaan Positif Lengkap}
Dalam contoh~\ref{001306} ditunjukkan bagaimana himpunan $\mathbf M_n$ semua matriks $n \times n$
atas bilangan kompleks dapat dijadikan aljabar-$C^*$. Sekarang kita memperumum contoh itu ke
aljabar $\M nA$ semua matriks $n \times n$ yang unsur-unsurnya berasal dari aljabar-$C^*$~$A$.

\begin{exam} Dalam contoh~\ref{000533a} dinyatakan bahwa jika $A$ merupakan aljabar-$C^*$, maka di bawah
operasi aljabar biasa, himpunan $\M nA$ semua matriks $n \times n$ yang unsur-unsurnya berasal dari
$A$ merupakan suatu aljabar. Aljabar ini dapat dijadikan aljabar-$*\,$ dengan mengambil transposisi
konjugat sebagai involusi. Artinya, definisikan
   \[ [ a_{ij}]^* := \bigl[ {a_{j\,i}}^*\bigr] \]
dengan $a_{ij} \in A$ untuk $1 \le i,j \le n$.

Misalkan $H$ ruang Hilbert. Untuk sementara, ruang itu tidak harus berdimensi tak hingga atau
separabel. Lambangkan dengan $H^n$ jumlah langsung rangkap-$n$-nya. Artinya,
   \[ H^n := \bigoplus_{k=1}^n H_k \]
dengan $H_k = H$ untuk $k = 1, \dots, n$. Misalkan $[T_{jk}] \in \mathbf M_n(\ofml B(H))$. Untuk
$\vc x = x_1 \oplus \dots \oplus x_n \in H^n$, definisikan
   \[ \vc T \colon H^n \sto H^n \colon \vc x \mapsto [T_{jk}] \vc x =
            \sum_{k=1}^n T_{1k}x_k \oplus \dots \oplus \sum_{k=1}^n T_{nk}x_k \,. \]
Maka $\vc T \in \ofml B(H^n)$. Lebih lanjut, pemetaan
   \[ \Psi \colon \mathbf M_n\bigl(\ofml B(H)\bigr) \sto \ofml B(H^n) \colon
             [T_{jk}] \mapsto \vc T \]
merupakan isomorfisme aljabar-$*\,$. Gunakan $\Psi$ untuk memindahkan norma operator pada $\ofml
B(H^n)$ ke~$\mathbf M_n\bigl(\ofml B(H)\bigr)$. Artinya, definisikan
   \[ \bignorm{\,[T_{jk}]\,} := \norm{\vc T}\,. \]
Ini menjadikan $\mathbf M_n\bigl(\ofml B(H)\bigr)$ aljabar-$C^*$ yang isomorfik dengan~$\ofml
B(H^n)$.

Sekarang andaikan $A$ aljabar-$C^*$ sebarang. Gunakan versi III Teorema
Gelfand--Naimark pada~\ref{002973}. Dengan teorema ini, kita dapat mengidentifikasi $A$ dengan
subaljabar-$C^*$ dari $\ofml B(H)$ untuk suatu ruang Hilbert $H$ dan, akibatnya,
 \index{M@$\M nA$!sebagai aljabar-$C^*$}%
 \index{C*@aljabar-$C^*$!M@$\M nA$ sebagai suatu}%
$\M nA$ dengan suatu subaljabar-$C^*$ dari $\mathbf M_n\bigl(\ofml B(H)\bigr)$. Dengan
identifikasi ini, $\M nA$ menjadi aljabar-$C^*$. (Perhatikan bahwa menurut akibat~\ref{00131104},
norma pada aljabar-$C^*$ bersifat tunggal; jadi jelas bahwa norma pada $\mathbf M_n(A)$
tidak bergantung pada cara khusus $A$ direpresentasikan sebagai subaljabar-$C^*$ dari
operator pada suatu ruang Hilbert.)
\end{exam}

\begin{exam} Misalkan $k \in \N$ dan $A = \mathbf M_k$. Maka untuk setiap $n \in \N$,
   \[ \mathbf M_n(A) = \mathbf M_n(\mathbf M_k) \cong \mathbf M_{nk}\,. \]
Isomorfismenya jelas: cukup hapus kurung siku bagian dalam. Sebagai contoh,
isomorfisme dari $\mathbf M_2(\mathbf M_2)$ ke $\mathbf M_4$ diberikan oleh
  \[ \begin{bmatrix} {} & {} \\
          \begin{bmatrix}
                          a_{11} & a_{12} \\
                          a_{21} & a_{22}
          \end{bmatrix}       &
          \begin{bmatrix} b_{11} & b_{12} \\
                          b_{21} & b_{22}
          \end{bmatrix}       \\*[25pt]
          \begin{bmatrix} c_{11} & c_{12} \\
                          c_{21} & c_{22}
          \end{bmatrix}       &
          \begin{bmatrix} d_{11} & d_{12} \\
                          d_{21} & d_{22}
          \end{bmatrix}       \\
                    {}  & {}
     \end{bmatrix}
    \quad \mapsto \quad
          \begin{bmatrix} a_{11} & a_{12} & b_{11} & b_{12} \\
                          a_{21} & a_{22} & b_{21} & b_{22} \\
                          c_{11} & c_{12} & d_{11} & d_{12} \\
                          c_{21} & c_{22} & d_{21} & d_{22}
          \end{bmatrix}
  \]
\end{exam}

\begin{defn} Suatu pemetaan $\phi\colon A \sto B$ di antara aljabar-$C^*$ disebut \df{positif} jika
pemetaan itu membawa elemen positif ke elemen positif; yakni, jika $\phi(a) \in B^+$ setiap kali
$a \in A^+$.
\end{defn}

\begin{exam} Setiap homomorfisme-$*\,$ di antara aljabar-$C^*$ bersifat positif. (Lihat
proposisi~\ref{0018201}.)
\end{exam}

\begin{prop} Setiap pemetaan linear positif di antara aljabar-$C^*$ bersifat terbatas.
\end{prop}

\begin{notn} Misalkan $A$ dan $B$ aljabar-$C^*$ serta $n \in \N$. Suatu pemetaan linear $\phi\colon A \sto
B$ menginduksi pemetaan linear $\phi^{(\,n)} := \phi \otimes \id{\mathbf M_n}\colon A \otimes
\mathbf M_n \sto B \otimes \mathbf M_n$. Seperti telah kita lihat, $A \otimes \mathbf M_n$ dapat
diidentifikasi dengan $\mathbf M_n(A)$. Jadi kita dapat memandang $\phi^{(\,n)}$ sebagai pemetaan dari
$\mathbf M_n(A)$ ke $\mathbf M_n(B)$ yang didefinisikan oleh $\phi^{(\,n)}\big(\,[a_{jk}]\,\bigr) =
\bigl[ \phi(a_{jk})\bigr]$.
\end{notn}

Mudah dilihat bahwa jika $\phi$ mempertahankan perkalian, maka setiap $\phi^{(\,n)}$ juga demikian,
dan jika $\phi$ mempertahankan involusi, maka setiap $\phi^{(\,n)}$ juga demikian.
Akan tetapi, kepositifan tidak selalu dipertahankan, sebagaimana ditunjukkan contoh~\ref{005824} berikut.

\begin{defn} Dalam $\mathbf M_n$,
 \index{standar!unit matriks}%
 \index{matriks!unit}%
 \index{unit!matriks}%
\df{unit matriks standar} adalah matriks-matriks $e^{jk}$ ($1 \le j,k \le n$) yang entri pada
baris ke-$j$ dan kolom ke-$k$ bernilai $1$, sedangkan semua entri lainnya
bernilai~$0$.
\end{defn}

\begin{exam}\label{005824} Misalkan $A = \mathbf M_2$. Maka $\phi\colon A \sto A\colon a \mapsto a^t$
(dengan $a^t$ transposisi~$a$) merupakan pemetaan positif. Pemetaan $\phi^{(2)}\colon
\mathbf M_2(A) \sto \mathbf M_2(A)$ tidak positif. Untuk melihatnya, misalkan $e^{11}$, $e^{12}$,
$e^{21}$, dan $e^{22}$ unit matriks standar untuk~$A$. Maka
  \[ \phi^{(2)}\left(\begin{bmatrix} e^{11} & e^{12} \\ e^{21} & e^{22} \end{bmatrix}\right)
    = \begin{bmatrix} \phi(e^{11}) & \phi(e^{12}) \\ \phi(e^{21}) & \phi(e^{22}) \end{bmatrix}
    = \begin{bmatrix} e^{11} & e^{21} \\ e^{12} & e^{22} \end{bmatrix}
    = \begin{bmatrix} 1 & 0 & 0 & 0 \\
                      0 & 0 & 1 & 0 \\
                      0 & 1 & 0 & 0 \\
                      0 & 0 & 0 & 1
      \end{bmatrix} \]
yang tidak positif.
\end{exam}

\begin{defn} Suatu pemetaan linear $\phi\colon A \sto B$ di antara aljabar-$C^*$ beridentitas disebut
 \index{npositif@$n$-positif}%
 \index{positif!$n$-}%
\df{$n$-positif} (untuk $n \in \N$) jika $\phi^{(\,n)}$ positif. Pemetaan itu disebut
 \index{positif!lengkap}%
 \index{lengkap!positif}%
\df{positif lengkap} jika bersifat $n$-positif untuk setiap $n \in \N$.
\end{defn}

\begin{exam} Setiap homomorfisme-$*\,$ di antara aljabar-$C^*$ bersifat positif lengkap.
\end{exam}

\begin{prop}\label{00582811} Misalkan $A$ aljabar-$C^*$ beridentitas dan $a \in A$. Maka $\norm a \le 1$ jika dan hanya jika
$\begin{bmatrix} \vc 1 & a \\ a^* & \vc 1 \end{bmatrix}$ positif dalam~$\mathbf M_2(A)$.
\end{prop}

\begin{prop} Misalkan $A$ aljabar-$C^*$ beridentitas dan $a$, $b \in A$. Maka $a^*a \le b$ jika
dan hanya jika $\begin{bmatrix} \vc 1 & a \\ a^* & b \end{bmatrix}$ positif dalam~$\mathbf
M_2(A)$.
\end{prop}

\begin{prop} Setiap pemetaan $2$-positif beridentitas di antara aljabar-$C^*$ beridentitas bersifat kontraktif.
\end{prop}

\begin{prop}[Ketaksamaan Kadison]\label{0058283} Jika $\phi\colon A \sto B$ merupakan pemetaan
$2$-positif beridentitas
 \index{ketaksamaan Kadison}%
 \index{ketaksamaan!Kadison}%
di antara aljabar-$C^*$ beridentitas, maka
   \[ \bigl(\phi(a)\bigr)^*\phi(a) \le \phi(a^*a) \]
untuk setiap $a \in A$.
\end{prop}

\begin{defn} Mudah dilihat bahwa jika $\phi\colon A \sto B$ merupakan pemetaan linear terbatas
di antara aljabar-$C^*$, maka $\phi^{(\,n)}$ terbatas untuk setiap~$n \in \N$. Jika
$\norm{\phi}_{\mathrm{cb}} := \sup\{\norm{\phi^{(\,n)}} \colon n \in \N\} < \infty$, maka
$\phi$ disebut
 \index{lengkap!terbatas}%
 \index{terbatas!lengkap}%
 \index{<unopn@$\norm{\hphantom{\phi}}_{\mathrm{cb}}$ (norma pemetaan terbatas lengkap)}%
\df{terbatas lengkap}.
\end{defn}

\begin{prop}\label{00582921} Jika $\phi\colon A \sto B$ merupakan pemetaan linear beridentitas dan positif lengkap
di antara aljabar-$C^*$ beridentitas, maka $\phi$ terbatas lengkap dan
   \[ \norm{\phi(\vc 1)} = \norm\phi = \norm\phi_{\mathrm{cb}}\,. \]
\end{prop}

\begin{proof} Lihat \cite{Paulsen:2002}. \ns \end{proof}

\begin{thm}[Teorema dilatasi Stinespring]\label{005831} Misalkan $A$ aljabar-$C^*$ beridentitas,
 \index{Stinespring!teorema dilatasi}%
 \index{dilatasi!teorema Stinespring}%
$H$ ruang Hilbert, dan $\phi\colon A \sto \ofml B(H)$ pemetaan linear beridentitas.
Maka $\phi$ positif lengkap jika dan hanya jika terdapat suatu ruang Hilbert
$H_0$, suatu isometri $V\colon H \sto H_0$, dan suatu representasi tak terdegenerasi $\vc
r\colon A \sto \ofml B(H_0)$ sedemikian sehingga $\phi(a) = V^*\vc r(a)V$ untuk setiap $a \in A$.
\end{thm}

\begin{proof} Lihat \cite{HigsonR:2000}, teorema 3.1.3, atau \cite{Paulsen:2002}, teorema 4.1. \ns
\end{proof}

\begin{notn} Misalkan $A$ aljabar-$C^*$, $f\colon A \sto \ofml B(H_0)$, dan
$g\colon A \sto \ofml B(H)$, dengan $H_0$ dan $H$ ruang Hilbert. Kita menulis
 \index{<relation@$g \lesssim f$ (relasi di antara pemetaan bernilai operator)}%
$g \lesssim f$ jika terdapat suatu isometri $V\colon H \sto H_0$ sedemikian sehingga $g(a) -
V^*f(a)V \in \ofml K(H)$ untuk setiap $a \in A$. Relasi $\lesssim$ merupakan praurutan,
tetapi bukan urutan parsial; yakni, relasi itu refleksif dan transitif, tetapi tidak
antisimetris.
\end{notn}

\begin{thm}[Voiculescu]\label{0058335} Misalkan $A$ aljabar-$C^*$ beridentitas yang separabel,
 \index{teorema Voiculescu}%
$\vc r\colon A \sto \ofml B(H_0)$ representasi tak terdegenerasi, dan
$\phi\colon A \sto \ofml B(H)$ pemetaan linear beridentitas dan positif lengkap (dengan $H$ dan $H_0$ ruang Hilbert). Jika $\phi(a) =
\vc 0$ setiap kali $\vc r(a) \in \ofml K(H_0)$, maka $\phi \lesssim \vc r$.
\end{thm}

\begin{proof} Buktinya panjang dan rumit, tetapi ``elementer''. Lihat \cite{HigsonR:2000},
bagian 3.5--3.6. \ns
\end{proof}

\begin{prop}\label{0058338} Misalkan $A$ aljabar-$C^*$ beridentitas yang separabel, $H$ ruang Hilbert,
dan $\tau_1$, $\tau_2 \colon A \sto \ofml Q(H)$ monomorfisme-$*\,$ beridentitas. Jika
$\tau_2$ terbelah, maka $\tau_1 + \tau_2$ ekuivalen secara uniter dengan~$\tau_1$.
\end{prop}

\begin{proof} Lihat \cite{HigsonR:2000}, teorema 3.4.7. \ns \end{proof}

\begin{cor}\label{0058339} Andaikan $A$ aljabar-$C^*$ beridentitas yang separabel, $H$ ruang
Hilbert, dan $\tau \colon A \sto \ofml Q(H)$ monomorfisme-$*\,$ beridentitas yang terbelah.
Maka $[\tau]$ merupakan identitas aditif dalam~$\ext A$.
\end{cor}

Perhatikan bahwa suatu monomorfisme-$*\,$ beridentitas $\tau\colon A \sto \ofml Q(H)$ dari aljabar-$C^*$
beridentitas yang separabel ke aljabar Calkin dari suatu ruang Hilbert $H$ bersifat semiterbelah jika dan
hanya jika $[\tau]$ mempunyai invers aditif dalam $\ext A$. Proposisi berikut mengatakan bahwa
hal ini terjadi jika dan hanya jika $\tau$ mempunyai suatu pengangkatan positif lengkap ke~$\ofml B(H)$.

\begin{prop}\label{0058411} Misalkan $A$ aljabar-$C^*$ beridentitas yang separabel dan $H$ ruang Hilbert.
Suatu monomorfisme-$*\,$ beridentitas $\tau\colon A \sto \ofml Q(H)$ semiterbelah jika dan hanya jika
terdapat suatu pemetaan linear beridentitas dan positif lengkap $\wt\tau\colon A \sto \ofml B(H)$ sedemikian sehingga
diagram
  \[
    \xy
      \dtriangle/<-`->`->/[\ofml B(H)`A`\ofml Q(H);\wt\tau`\pi`\tau]
    \endxy
   \]
komutatif.
\end{prop}

\begin{proof} Lihat \cite{HigsonR:2000}, teorema 3.1.5.  \ns \end{proof}

\begin{defn} Suatu aljabar-$C^*$ $A$ disebut
 \index{nuklir}%
 \index{C*@aljabar-$C^*$!nuklir}%
\df{nuklir} jika untuk setiap aljabar-$C^*$ $B$ terdapat suatu norma-$C^*$ tunggal pada hasil kali
tensor aljabaris $A \odot B$.
\end{defn}

\begin{exam}\label{0058513} Setiap aljabar-$C^*$ berdimensi hingga bersifat nuklir.
 \index{nuklir!aljabar berdimensi hingga bersifat}%
\end{exam}

\begin{proof} Lihat \cite{Murphy:1990}, teorema 6.3.9.  \ns \end{proof}

\begin{exam}\label{0058514} Aljabar-$C^*$ $\mathbf M_n$
 \index{M@$\mathbf M_n$!sebagai aljabar-$C^*$ nuklir}%
 \index{nuklir!$\mathbf M_n$ bersifat}%
bersifat nuklir.
\end{exam}
\begin{proof} Lihat \cite{Blackadar:2006}, II.9.4.2.  \ns \end{proof}

\begin{exam}\label{0058515} Jika $H$ ruang Hilbert,
 \index{K@$\ofml K(H)$!sebagai aljabar-$C^*$ nuklir}%
 \index{nuklir!$\ofml K(H)$ bersifat}%
maka aljabar-$C^*$ $\ofml K(H)$ dari operator kompak bersifat nuklir.
\end{exam}
\begin{proof} Lihat \cite{Murphy:1990}, contoh 6.3.2.  \ns \end{proof}

\begin{exam}\label{0058517} Jika $X$ ruang Hausdorff kompak, maka aljabar-$C^*$
 \index{kontinu@$\fml C(X)$!sebagai aljabar-$C^*$ nuklir}%
 \index{nuklir!$\fml C(X)$ bersifat}%
$\fml C(X)$ bersifat nuklir.
\end{exam}

\begin{proof} Lihat \cite{Paulsen:2002}, proposisi 12.9.  \ns \end{proof}

\begin{exam}\label{0058519} Setiap aljabar-$C^*$ komutatif bersifat nuklir.
 \index{nuklir!aljabar-$C^*$ komutatif bersifat}%
\end{exam}

\begin{proof} Lihat \cite{Fillmore:1996}, teorema 7.4.1, atau \cite{Wegge-Olsen:1993}, teorema
T.6.20. \ns
\end{proof}

Kesimpulan hasil berikut dikenal sebagai
 \index{lengkap!positif!sifat pengangkatan}%
\emph{sifat pengangkatan positif lengkap}.

\begin{prop}\label{0058541} Misalkan $A$ aljabar-$C^*$ beridentitas, separabel, dan nuklir, serta $J$ suatu
ideal separabel dalam aljabar-$C^*$ beridentitas~$B$. Maka untuk setiap pemetaan positif lengkap
$\phi\colon A \sto B/J$, terdapat suatu pemetaan positif lengkap $\wt\phi\colon A \sto B$
yang membuat diagram
  \[
    \xy
      \dtriangle/<-`->`->/[B`A`B/J;\wt\phi``\phi]
    \endxy
   \]
komutatif.
\end{prop}

\begin{proof} Lihat \cite{HigsonR:2000}, teorema 3.3.6.  \ns \end{proof}

\begin{cor}\label{0058543} Jika $A$ merupakan aljabar-$C^*$ beridentitas, separabel, dan nuklir,
 \index{ext@$\ext A$, $\ext X$!merupakan grup Abelian}%
maka $\ext A$ merupakan grup Abelian.
\end{cor}







\endinput
